\chapter{Índice Tentativo de Proyecto de Tesis}
\label{chap:indice}

\begin{tabular}{p{17cm}}
\noindent 
\textbf{RESUMEN} \\
\textbf{CAPÍTULO I} \\
\textbf{1. INTRODUCCIÓN} \\
\begin{enumerate}[leftmargin=*]
    \item[1.1] Motivación y Contexto
    \begin{itemize}
        \item Crisis de calidad en aplicaciones Ionic
        \item Costos ocultos del testing manual
    \end{itemize}
    
    \item[1.2] Planteamiento del Problema
    \begin{enumerate}
        \item[1.2.1] Problema General: Ineficiencia en generación/mantenimiento de pruebas unitarias
        \item[1.2.2] Problemas Específicos:
        \begin{itemize}
            \item Falta de priorización inteligente
            \item Pruebas desactualizadas
            \item Dependencia del conocimiento tribal
        \end{itemize}
    \end{enumerate}
    
    \item[1.3] Objetivos
    \begin{enumerate}
        \item[1.3.1] Objetivo General: Sistema de generación adaptativa con IA
        \item[1.3.2] Objetivos Específicos:
        \begin{itemize}
            \item Módulo Python para análisis estático
            \item Algoritmo de clustering por niveles
            \item Integración CI/CD
        \end{itemize}
    \end{enumerate}
    
    \item[1.4] Contribuciones Originales
    \begin{itemize}
        \item Modelo CLUSTest-Ionic
        \item Patrón Test-As-Code
        \item Dataset abierto
    \end{itemize}
    
    \item[1.5] Organización de la Tesis
\end{enumerate}

\textbf{CAPÍTULO II} \\
\textbf{2. MARCO TEÓRICO} \\
\begin{enumerate}[leftmargin=*]
    \item[2.1] Fundamentos de Testing en Ionic
    \begin{itemize}
        \item Arquitectura Angular/Ionic
        \item Pirámide de pruebas
    \end{itemize}
    
    \item[2.2] Automatización Inteligente
    \begin{itemize}
        \item LLMs para generación de código
        \item Fine-tuning para dominios específicos
    \end{itemize}
    
    \item[2.3] Métricas de Calidad
    \begin{itemize}
        \item Complejidad ciclomática adaptada
        \item Cobertura efectiva
    \end{itemize}
\end{enumerate}

\textbf{CAPÍTULO III} \\
\textbf{3. TRABAJOS RELACIONADOS} \\
\begin{itemize}
    \item Herramientas existentes (Jest, Spectator)
    \item Limitaciones en soluciones actuales
    \item Brecha de investigación identificada
\end{itemize}
\end{tabular}

\begin{tabular}{p{17cm}}
\textbf{CAPÍTULO IV} \\
\textbf{4. PROPUESTA} \\
\begin{enumerate}[leftmargin=*]
    \item[4.1] Arquitectura del Sistema
    \begin{itemize}
        \item Diagrama de componentes
        \item Flujo de datos
    \end{itemize}
    
    \item[4.2] Implementación
    \begin{enumerate}
        \item[4.2.1] Módulo de Análisis
        \begin{itemize}
            \item AST Parser
            \item Calculador de complejidad
        \end{itemize}
        
        \item[4.2.2] Generador Adaptativo
        \begin{itemize}
            \item Integración DeepSeek
            \item Templates por niveles
        \end{itemize}
        
        \item[4.2.3] CI/CD Integration
        \begin{itemize}
            \item GitHub Actions
            \item Reportes automáticos
        \end{itemize}
    \end{enumerate}
    
    \item[4.3] Innovaciones Clave
    \begin{itemize}
        \item Clasificación por criticidad
        \item Auto-corrección de pruebas
    \end{itemize}
\end{enumerate}

\textbf{CAPÍTULO V} \\
\textbf{5. VALIDACIÓN EXPERIMENTAL} \\
\begin{enumerate}[leftmargin=*]
    \item[5.1] Metodología
    \begin{itemize}
        \item Proyectos de evaluación
        \item Métricas comparativas
    \end{itemize}
    
    \item[5.2] Resultados Cuantitativos
    \begin{itemize}
        \item Tiempos de generación
        \item Cobertura alcanzada
    \end{itemize}
    
    \item[5.3] Resultados Cualitativos
    \begin{itemize}
        \item Feedback de desarrolladores
        \item Casos de estudio
    \end{itemize}
\end{enumerate}

\textbf{CAPÍTULO VI} \\
\textbf{6. CONCLUSIONES} \\
\begin{enumerate}[leftmargin=*]
    \item[6.1] Hallazgos Principales
    \item[6.2] Limitaciones
    \item[6.3] Trabajos Futuros
    \begin{itemize}
        \item Extensión a pruebas E2E
        \item Soporte para React/Vue
    \end{itemize}
\end{enumerate}

\textbf{ANEXOS} \\
\begin{itemize}
    \item Manual de instalación
    \item Ejemplos de código
\end{itemize}

\textbf{BIBLIOGRAFÍA} \\
\textbf{GLOSARIO} \\
\end{tabular}
